%! Unit 2: Functions and Their Graphs — Summative Assessment — Unit Test (blank)
%! Absolute Value & Radical Functions. A printed parallel of the teacher's
%! DeltaMath "Absolute Value & Radicals Test" (27 items): the same problem type
%! in the same order, with new functions and numbers. The DeltaMath items that
%! ask a student to drag a graph into place are given here as a table of
%! landmark points to fill (no freehand graphing on paper).
%! It is NEVER merged into a packet; distribute it separately. Its key must
%! paginate identically to it, so every worked solution lives in a
%! \begin{work} block authored byte-identically in both files.
\documentclass[10pt]{article}
\usepackage{precalculus-article}
\usepackage{precalculus-boxes}
\usepackage{grading}
\gradingassignment{precalc-u02-actual-test}
\gradingcourse{precalculus}

% Part-divider strip
\newcommand{\parthead}[1]{%
  \vspace{6pt}\noindent
  \begin{headlinebox}{plum}{\color{white}\bfseries #1}\end{headlinebox}%
  \vspace{4pt}}

% Question number
\newcommand{\qn}[1]{\par\noindent\textbf{#1.}\ }

% Coordinate grid: \gridaxes{xmin}{xmax}{ymin}{ymax} (integers; the window
% must contain the origin). A light grid line at EVERY unit and a label at EVERY
% nonzero integer, so a student can read any lattice point exactly. (TikZ's
% default grid step is 1cm, not 1 unit: always pass step=1.) Labels sit on a
% white patch so grid lines never cross the digits; draw curves after the call
% so they stay on top.
\newcommand{\gridaxes}[4]{%
  \draw[linegray!55, thin, step=1] (#1,#3) grid (#2,#4);
  \draw[thick, ->, >=stealth] (#1-0.4,0) -- (#2+0.8,0) node[right, inner sep=1pt]{\small $x$};
  \draw[thick, ->, >=stealth] (0,#3-0.4) -- (0,#4+0.8) node[above, inner sep=1pt]{\small $y$};
  \foreach \t in {#1,...,#2}{\ifnum\t=0\else
    \draw[thick] (\t,0.15) -- (\t,-0.15)
      node[below, inner sep=0.6pt, fill=white, font=\scriptsize]{$\t$};\fi}
  \foreach \t in {#3,...,#4}{\ifnum\t=0\else
    \draw[thick] (0.15,\t) -- (-0.15,\t)
      node[left, inner sep=0.6pt, fill=white, font=\scriptsize]{$\t$};\fi}}

% A plotted point: a small filled dot centred exactly on the lattice point
\newcommand{\gpt}[2][plum]{\fill[#1] #2 circle (2.3pt)}

\begin{document}

\pageheader{Unit 2: Functions \& Their Graphs}{Unit Test: Absolute Value \& Radical Functions}

\vspace{0.06in}
\noindent\textbf{Instructions:} Show all work in the space provided. Write each final answer
in its box. Only what is inside a box is graded. Write the letter of the single best answer
in its Letter box. Write intervals in interval notation. No notes. \textbf{70 points total.}

% ======================================================================
\boxguard[6]
\parthead{Part A --- Transformations (12 pts)}

\noindent
\begin{minipage}[t]{0.5\textwidth}
\qn{1} The solid graph is $f(x) = \sqrt{x}$. The dashed graph is a transformation of $f$.
Fill in the boxes so the equation matches the dashed graph.

\vspace{6pt}
\noindent\hspace*{1em}$g(x) = $ \gradedblank[id=p1a, type=short, points=1, answer=1, check=numeric]{0.9cm}{1.1cm}
$f\big($\gradedblank[id=p1b, type=short, points=1, answer=-1, check=numeric]{0.9cm}{1.1cm}$\,x\ +$
\gradedblank[id=p1c, type=short, points=1, answer=2, check=numeric]{0.9cm}{1.1cm}$\big)\ +$
\gradedblank[id=p1d, type=short, points=1, answer=-3, check=numeric]{0.9cm}{1.1cm}
\end{minipage}%
\hfill
\begin{minipage}[t]{0.46\textwidth}
\vspace{0pt}\centering
\begin{tikzpicture}[x=0.44cm, y=0.44cm]
  \gridaxes{-8}{9}{-4}{4}
  \draw[plum, very thick, domain=0:9, samples=120, smooth] plot (\x, {sqrt(\x)});
  \draw[redacc, very thick, dashed, domain=-7:2, samples=200, smooth] plot (\x, {sqrt(2-\x)-3});
  \gpt[redacc]{(2,-3)};
  \gpt{(0,0)};
\end{tikzpicture}
\end{minipage}

\vspace{8pt}
\boxguard[14]
\qn{2} The graph of $y = 2\sqrt{x + 3}$ is the graph of the parent $y = \sqrt{x}$ after
transformations. Complete the table: where does each landmark point of the parent land?

\vspace{4pt}
\begin{center}
\renewcommand{\arraystretch}{1.5}
\begin{tabular}{|l|c|c|c|}
\hline
Parent point on $y = \sqrt{x}$ & $(0,\,0)$ & $(1,\,1)$ & $(4,\,2)$ \\ \hline
Lands at on $y = 2\sqrt{x+3}$ &
  \gradedblank[id=p2a, type=short, points=1, answer={(-3, 0)}, check=text, display={$(-3,\,0)$}]{2cm}{1.1cm} &
  \gradedblank[id=p2b, type=short, points=1, answer={(-2, 2)}, check=text, display={$(-2,\,2)$}]{2cm}{1.1cm} &
  \gradedblank[id=p2c, type=short, points=1, answer={(1, 4)}, check=text, display={$(1,\,4)$}]{2cm}{1.1cm} \\ \hline
\end{tabular}
\end{center}

\vspace{4pt}
\qn{3} The graph of $y = \left|\tfrac{1}{3}x\right| - 4$ is the graph of the parent
$y = |x|$ after transformations. Complete the table.

\vspace{4pt}
\begin{center}
\renewcommand{\arraystretch}{1.5}
\begin{tabular}{|l|c|c|c|}
\hline
Parent point on $y = |x|$ & $(-1,\,1)$ & $(0,\,0)$ & $(1,\,1)$ \\ \hline
Lands at on $y = \left|\tfrac{1}{3}x\right| - 4$ &
  \gradedblank[id=p3a, type=short, points=1, answer={(-3, -3)}, check=text, display={$(-3,\,-3)$}]{2cm}{1.1cm} &
  \gradedblank[id=p3b, type=short, points=1, answer={(0, -4)}, check=text, display={$(0,\,-4)$}]{2cm}{1.1cm} &
  \gradedblank[id=p3c, type=short, points=1, answer={(3, -3)}, check=text, display={$(3,\,-3)$}]{2cm}{1.1cm} \\ \hline
\end{tabular}
\end{center}

\vspace{6pt}
\boxguard[16]
\noindent
\begin{minipage}[t]{0.5\textwidth}
\qn{4} The solid graph is $y = |x|$. Which function is the dashed graph?

\vspace{6pt}
\hspace*{1em}(A) $y = -|x - 2| + 1$\par
\hspace*{1em}(B) $y = -|x + 2| + 1$\par
\hspace*{1em}(C) $y = |x - 2| + 1$\par
\hspace*{1em}(D) $y = -|x - 1| + 2$\par
\vspace{6pt}
\hspace*{1em}\textit{Letter:}~\gradedblank[id=p4, type=choice, points=2, answer=A]{1.1cm}{13pt}
\end{minipage}%
\hfill
\begin{minipage}[t]{0.46\textwidth}
\vspace{0pt}\centering
\begin{tikzpicture}[x=0.4cm, y=0.4cm]
  \gridaxes{-6}{8}{-6}{6}
  \draw[plum, very thick, <->] (-5.5,5.5) -- (0,0) -- (5.5,5.5);
  \draw[redacc, very thick, dashed, <->] (-4.5,-5.5) -- (2,1) -- (7.5,-4.5);
  \gpt{(0,0)};
  \gpt[redacc]{(2,1)};
\end{tikzpicture}
\end{minipage}

% ======================================================================
\boxguard[30]
\parthead{Part B --- Symmetry, One-to-One, and Inverses (14 pts)}

\qn{5} Fill in each table of values. Then write whether the function is \emph{even},
\emph{odd}, or \emph{neither}.

\vspace{4pt}
\noindent
\begin{minipage}[t]{0.32\textwidth}
\centering $f(x) = 2x^3 - x$\par\vspace{3pt}
\renewcommand{\arraystretch}{1.35}
\begin{tabular}{|c|c|}
\hline $x$ & $y$ \\ \hline
$-2$ & \gradedblank[id=p5a1, type=short, points=0.2, answer=-14, check=numeric]{1.1cm}{1.1cm} \\ \hline
$-1$ & \gradedblank[id=p5a2, type=short, points=0.2, answer=-1, check=numeric]{1.1cm}{1.1cm} \\ \hline
$0$  & \gradedblank[id=p5a3, type=short, points=0.2, answer=0, check=numeric]{1.1cm}{1.1cm} \\ \hline
$1$  & \gradedblank[id=p5a4, type=short, points=0.2, answer=1, check=numeric]{1.1cm}{1.1cm} \\ \hline
$2$  & \gradedblank[id=p5a5, type=short, points=0.2, answer=14, check=numeric]{1.1cm}{1.1cm} \\ \hline
\end{tabular}\par\vspace{11pt}
\gradedblank[id=p5a, type=short, points=1, answer=odd, check=text, display={Odd}]{2cm}{1.1cm}
\end{minipage}%
\hfill
\begin{minipage}[t]{0.32\textwidth}
\centering $g(x) = |x| + 3$\par\vspace{3pt}
\renewcommand{\arraystretch}{1.35}
\begin{tabular}{|c|c|}
\hline $x$ & $y$ \\ \hline
$-2$ & \gradedblank[id=p5b1, type=short, points=0.2, answer=5, check=numeric]{1.1cm}{1.1cm} \\ \hline
$-1$ & \gradedblank[id=p5b2, type=short, points=0.2, answer=4, check=numeric]{1.1cm}{1.1cm} \\ \hline
$0$  & \gradedblank[id=p5b3, type=short, points=0.2, answer=3, check=numeric]{1.1cm}{1.1cm} \\ \hline
$1$  & \gradedblank[id=p5b4, type=short, points=0.2, answer=4, check=numeric]{1.1cm}{1.1cm} \\ \hline
$2$  & \gradedblank[id=p5b5, type=short, points=0.2, answer=5, check=numeric]{1.1cm}{1.1cm} \\ \hline
\end{tabular}\par\vspace{11pt}
\gradedblank[id=p5b, type=short, points=1, answer=even, check=text, display={Even}]{2cm}{1.1cm}
\end{minipage}%
\hfill
\begin{minipage}[t]{0.32\textwidth}
\centering $h(x) = |x + 1|$\par\vspace{3pt}
\renewcommand{\arraystretch}{1.35}
\begin{tabular}{|c|c|}
\hline $x$ & $y$ \\ \hline
$-2$ & \gradedblank[id=p5c1, type=short, points=0.2, answer=1, check=numeric]{1.1cm}{1.1cm} \\ \hline
$-1$ & \gradedblank[id=p5c2, type=short, points=0.2, answer=0, check=numeric]{1.1cm}{1.1cm} \\ \hline
$0$  & \gradedblank[id=p5c3, type=short, points=0.2, answer=1, check=numeric]{1.1cm}{1.1cm} \\ \hline
$1$  & \gradedblank[id=p5c4, type=short, points=0.2, answer=2, check=numeric]{1.1cm}{1.1cm} \\ \hline
$2$  & \gradedblank[id=p5c5, type=short, points=0.2, answer=3, check=numeric]{1.1cm}{1.1cm} \\ \hline
\end{tabular}\par\vspace{11pt}
\gradedblank[id=p5c, type=short, points=1, answer=neither, check=text, display={Neither}]{2cm}{1.1cm}
\end{minipage}

\vspace{12pt}
\boxguard[18]
\noindent For 6--8, choose from the answers under each graph.

\vspace{4pt}
\noindent
\begin{minipage}[t]{0.32\textwidth}
\qn{6} Is the function even, odd, or neither?\par
\vspace{3pt}\centering
\begin{tikzpicture}[x=0.46cm, y=0.46cm]
  \gridaxes{-4}{4}{-3}{3}
  \draw[plum, very thick, domain=-3.3:3.3, samples=120, smooth, <->] plot (\x, {0.1*(\x)^4 - (\x)^2 + 1});
\end{tikzpicture}\par
\raggedright\small
(A) even \quad (B) odd \quad (C) neither\par\vspace{4pt}
\textit{Letter:}~\gradedblank[id=p6, type=choice, points=2, answer=A]{1.1cm}{13pt}
\end{minipage}%
\hfill
\begin{minipage}[t]{0.32\textwidth}
\qn{7} Use the vertical and horizontal line tests. The graph shows
\par\vspace{3pt}\centering
\begin{tikzpicture}[x=0.46cm, y=0.46cm]
  \gridaxes{-4}{5}{-4}{3}
  \draw[plum, very thick, domain=-2.4:4.4, samples=120, smooth, <->] plot (\x, {0.5*(\x-1)^2 - 3});
\end{tikzpicture}\par
\raggedright\small
(A) a one-to-one function\par
(B) a function, not one-to-one\par
(C) not a function\par\vspace{4pt}
\textit{Letter:}~\gradedblank[id=p7, type=choice, points=2, answer=B]{1.1cm}{13pt}
\end{minipage}%
\hfill
\begin{minipage}[t]{0.32\textwidth}
\qn{8} Use the vertical and horizontal line tests. The graph shows
\par\vspace{3pt}\centering
\begin{tikzpicture}[x=0.46cm, y=0.46cm]
  \gridaxes{-4}{4}{-4}{4}
  \draw[plum, very thick, domain=-3.2:3.2, samples=120, smooth, <->] plot ({0.5*(\x)^2 - 3}, \x);
\end{tikzpicture}\par
\raggedright\small
(A) a one-to-one function\par
(B) a function, not one-to-one\par
(C) not a function\par\vspace{4pt}
\textit{Letter:}~\gradedblank[id=p8, type=choice, points=2, answer=C]{1.1cm}{13pt}
\end{minipage}

\vspace{12pt}
\boxguard[12]
\qn{9} Let $f(x) = \sqrt[3]{2x - 5}$. Which function is $f^{-1}(x)$?

\begin{work}
          y &= \sqrt[3]{2x - 5} \\
          x &= \sqrt[3]{2y - 5} \\
        x^3 &= 2y - 5 \\
          y &= \frac{x^3 + 5}{2}
\end{work}

\hspace*{1em}(A) $\dfrac{x^3 + 5}{2}$ \qquad (B) $\dfrac{x^3 - 5}{2}$ \qquad
(C) $\dfrac{x^3}{2} + 5$ \qquad (D) $\dfrac{1}{\sqrt[3]{2x - 5}}$
\hfill \textit{Letter:}~\gradedblank[id=p9, type=choice, points=2, answer=A]{1.1cm}{13pt}

% ======================================================================
\boxguard[14]
\parthead{Part C --- Reading Characteristics (21 pts)}

\noindent
\begin{minipage}[t]{0.5\textwidth}
\qn{10} The graph of $f$ continues in the direction of its arrows. Its turning points are
marked.

\vspace{8pt}
Increasing on:\par\vspace{2pt}
\gradedblank[id=p10a, type=short, points=1, answer={(-inf, -3) U (2, inf)}, check=text, display={$(-\infty,-3)\cup(2,\infty)$}]{4.6cm}{1.1cm}\par
\vspace{14pt}
Decreasing on:\par\vspace{2pt}
\gradedblank[id=p10b, type=short, points=1, answer={(-3, 2)}, check=text, display={$(-3,\,2)$}]{4.6cm}{1.1cm}
\end{minipage}%
\hfill
\begin{minipage}[t]{0.46\textwidth}
\vspace{0pt}\centering
\begin{tikzpicture}[x=0.5cm, y=0.5cm]
  \gridaxes{-6}{5}{-4}{4}
  \draw[plum, very thick, domain=-5.3:4.2, samples=150, smooth, <->]
    plot (\x, {0.24*((\x)^3/3 + (\x)^2/2 - 6*\x) - 0.24});
  \gpt{(-3,3)} node[above, inner sep=2pt, fill=white, font=\scriptsize]{$(-3,\,3)$};
  \gpt{(2,-2)} node[below, inner sep=2pt, fill=white, font=\scriptsize]{$(2,\,-2)$};
\end{tikzpicture}
\end{minipage}

\vspace{12pt}
\boxguard[18]
\noindent
\begin{minipage}[t]{0.48\textwidth}
\qn{11} State the domain of the function graphed.
\par\vspace{3pt}\centering
\begin{tikzpicture}[x=0.48cm, y=0.48cm]
  \gridaxes{-2}{10}{-2}{4}
  \draw[plum, very thick, domain=1:10, samples=150, smooth, ->] plot (\x, {3 - sqrt(\x-1)});
  \gpt{(1,3)};
\end{tikzpicture}\par\vspace{4pt}
\raggedright Domain: \gradedblank[id=p11, type=short, points=2, answer={[1, inf)}, check=text, display={$[1,\infty)$}]{3cm}{1.1cm}
\end{minipage}%
\hfill
\begin{minipage}[t]{0.48\textwidth}
\qn{12} State the range of the function graphed.
\par\vspace{3pt}\centering
\begin{tikzpicture}[x=0.48cm, y=0.48cm]
  \gridaxes{-8}{4}{-2}{3}
  \draw[plum, very thick, domain=-6:3, samples=150, smooth, <-] plot (\x, {sqrt(3-\x) - 1});
  \gpt{(3,-1)};
\end{tikzpicture}\par\vspace{4pt}
\raggedright Range: \gradedblank[id=p12, type=short, points=2, answer={[-1, inf)}, check=text, display={$[-1,\infty)$}]{3cm}{1.1cm}
\end{minipage}

\vspace{12pt}
\qn{13} Find the domain of $f(x) = \sqrt{10 - 2x} + 1$ algebraically.

\begin{work}
  10 - 2x &\ge 0 \\
      -2x &\ge -10 \\
        x &\le 5
\end{work}

\hspace*{1em}Domain: \gradedblank[id=p13, type=short, points=2, answer={(-inf, 5]}, check=text, display={$(-\infty,\,5]$}]{3cm}{1.1cm}

\vspace{12pt}
\boxguard[18]
\noindent
\begin{minipage}[t]{0.32\textwidth}
\qn{14} Is the graph concave up or concave down?
\par\vspace{3pt}\centering
\begin{tikzpicture}[x=0.42cm, y=0.42cm]
  \gridaxes{-6}{5}{-4}{4}
  \draw[plum, very thick, domain=-5:4, samples=150, smooth, ->] plot (\x, {2*sqrt(\x+5) - 3});
  \gpt{(-5,-3)};
\end{tikzpicture}\par
\raggedright\small
(A) concave up \quad (B) concave down\par\vspace{4pt}
\textit{Letter:}~\gradedblank[id=p14, type=choice, points=2, answer=B]{1.1cm}{13pt}
\end{minipage}%
\hfill
\begin{minipage}[t]{0.32\textwidth}
\qn{15} Is the graph concave up or concave down?
\par\vspace{3pt}\centering
\begin{tikzpicture}[x=0.42cm, y=0.42cm]
  \gridaxes{-2}{5}{-4}{4}
  \draw[plum, very thick, domain=-1:4.2, samples=150, smooth, ->] plot (\x, {0.25*(\x+1)^2 - 3});
  \gpt{(-1,-3)};
\end{tikzpicture}\par
\raggedright\small
(A) concave up \quad (B) concave down\par\vspace{4pt}
\textit{Letter:}~\gradedblank[id=p15, type=choice, points=2, answer=A]{1.1cm}{13pt}
\end{minipage}%
\hfill
\begin{minipage}[t]{0.32\textwidth}
\qn{16} Find the point(s) of inflection.
\par\vspace{3pt}\centering
\begin{tikzpicture}[x=0.42cm, y=0.42cm]
  \gridaxes{-4}{6}{-2}{6}
  \draw[plum, very thick, domain=-3.2:5.2, samples=150, smooth, <->] plot (\x, {0.125*(\x-1)^3 - 1.5*(\x-1) + 2});
\end{tikzpicture}\par
\raggedright\small
How many? \gradedblank[id=p16a, type=short, points=1, answer=1, check=numeric]{1.1cm}{1.1cm}\par\vspace{6pt}
At: \gradedblank[id=p16b, type=short, points=1, answer={(1, 2)}, check=text, display={$(1,\,2)$}]{2cm}{1.1cm}
\end{minipage}

\vspace{12pt}
\boxguard[22]
\noindent
\begin{minipage}[t]{0.56\textwidth}
\qn{17} The function $f$ is graphed at the right.

\vspace{4pt}
\begin{enumerate}[label=(\alph*), leftmargin=1.6em, itemsep=6pt]
  \item Domain of $f$: \gradedblank[id=p17a, type=short, points=1, answer={[-3, inf)}, check=text, display={$[-3,\infty)$}]{2.6cm}{1.1cm}
  \item Range of $f$: \gradedblank[id=p17b, type=short, points=1, answer={(-inf, 4]}, check=text, display={$(-\infty,\,4]$}]{2.6cm}{1.1cm}
  \item Which sentence is true?\par
    (A) $f$ is invertible because it passes the horizontal line test.\par
    (B) $f$ is invertible because it passes the vertical line test.\par
    (C) $f$ is not invertible because it is not one-to-one.\par
    \vspace{3pt}\textit{Letter:}~\gradedblank[id=p17c, type=choice, points=1, answer=A]{1.1cm}{13pt}
  \item Domain of $f^{-1}$: \gradedblank[id=p17d, type=short, points=1, answer={(-inf, 4]}, check=text, display={$(-\infty,\,4]$}]{2.6cm}{1.1cm}
  \item Range of $f^{-1}$: \gradedblank[id=p17e, type=short, points=1, answer={[-3, inf)}, check=text, display={$[-3,\infty)$}]{2.6cm}{1.1cm}
\end{enumerate}
\end{minipage}%
\hfill
\begin{minipage}[t]{0.4\textwidth}
\vspace{0pt}\centering
\begin{tikzpicture}[x=0.46cm, y=0.46cm]
  \gridaxes{-4}{9}{-2}{5}
  \draw[plum, very thick, domain=-3:9, samples=150, smooth, ->] plot (\x, {4 - sqrt(\x+3)});
  \gpt{(-3,4)};
\end{tikzpicture}
\end{minipage}

\vspace{12pt}
\boxguard[18]
\noindent
\begin{minipage}[t]{0.5\textwidth}
\qn{18} Which type of function is graphed?

\vspace{6pt}
\hspace*{1em}(A) exponential \qquad (D) square root\par
\hspace*{1em}(B) absolute value \quad\ (E) linear\par
\hspace*{1em}(C) cube root\par
\vspace{6pt}
\hspace*{1em}\textit{Letter:}~\gradedblank[id=p18, type=choice, points=2, answer=C]{1.1cm}{13pt}
\end{minipage}%
\hfill
\begin{minipage}[t]{0.46\textwidth}
\vspace{0pt}\centering
\begin{tikzpicture}[x=0.42cm, y=0.42cm]
  \gridaxes{-8}{8}{-5}{5}
  \draw[plum, very thick, domain=-7.5:7.5, samples=200, smooth, <->]
    plot (\x, {2*sign(\x)*pow(abs(\x),1/3)});
\end{tikzpicture}
\end{minipage}

% ======================================================================
\boxguard[16]
\parthead{Part D --- Average Rate of Change and Secant Lines (6 pts)}

Let $h(x) = \sqrt[3]{x - 3} - 1$.

\begin{work}
                      h(2) &= \sqrt[3]{-1} - 1 \\
                           &= -2 \\
                     h(11) &= \sqrt[3]{8} - 1 \\
                           &= 1 \\
  \frac{h(11) - h(2)}{11 - 2} &= \frac{1 - (-2)}{9} \\
                           &= \frac{1}{3} \\
                    y - (-2) &= \tfrac{1}{3}(x - 2) \\
                           y &= \tfrac{1}{3}x - \tfrac{8}{3}
\end{work}

\qn{19} Find the average rate of change of $h$ from $x = 2$ to $x = 11$. Write your answer
as a simplified fraction. \hfill
\gradedblank[id=p19, type=short, points=3, answer=1/3]{2cm}{1.1cm}

\vspace{10pt}
\qn{20} Find the secant line of $h$ from $x = 2$ to $x = 11$, in slope-intercept form.
\hfill $y = $ \gradedblank[id=p20, type=short, points=3, answer={x/3-8/3}, display={$\tfrac{1}{3}x - \tfrac{8}{3}$}]{3.4cm}{1.1cm}

% ======================================================================
\boxguard[16]
\parthead{Part E --- One Function, Every Characteristic (17 pts)}

Problems 21--27 are all about the function
\[
  g(x) = 3\,|x - 1| - 6 .
\]

\qn{21} Describe the transformations that take $y = |x|$ to $g$.

\rubricbox{\gradedblank[id=p21, type=rubric, points=3, display={Shift right $1$, stretch vertically by a factor of $3$, then shift down $6$.}]{\linewidth}{10pt}}

\vspace{8pt}
\qn{22} The domain of $g$ is
\gradedblank[id=p22a, type=short, points=1, answer={(-inf, inf)}, check=text, display={$(-\infty,\infty)$}]{2.8cm}{1.1cm},
and the range is
\gradedblank[id=p22b, type=short, points=1, answer={[-6, inf)}, check=text, display={$[-6,\infty)$}]{2.8cm}{1.1cm}.

\vspace{10pt}
\qn{23} $g$ is increasing on
\gradedblank[id=p23a, type=short, points=1, answer={(1, inf)}, check=text, display={$(1,\infty)$}]{2.8cm}{1.1cm}
and decreasing on
\gradedblank[id=p23b, type=short, points=1, answer={(-inf, 1)}, check=text, display={$(-\infty,\,1)$}]{2.8cm}{1.1cm}.

\vspace{6pt}
\boxguard[12]
\qn{24} Find every intercept of $g$.

\begin{work}
  3|x - 1| - 6 &= 0 \\
      |x - 1| &= 2 \\
        x - 1 &= 2 \ \text{or}\ x - 1 = -2 \\
            x &= 3 \ \text{or}\ x = -1 \\
         g(0) &= 3|0 - 1| - 6 \\
              &= -3
\end{work}

\hspace*{1em}$x$-intercept(s): $x = $
\gradedblank[id=p24a, type=short, points=2, answer={-1, 3}, check=set, display={$-1,\ 3$}]{2.6cm}{1.1cm}
\qquad $y$-intercept:
\gradedblank[id=p24b, type=short, points=1, answer={(0, -3)}, check=text, display={$(0,\,-3)$}]{2.6cm}{1.1cm}

\vspace{12pt}
\boxguard[12]
\noindent For 25 and 26, choose from these four statements.

\vspace{3pt}
\hspace*{1em}(A) $\displaystyle\lim_{x\to\infty} g(x) = \infty$ \qquad
(B) $\displaystyle\lim_{x\to-\infty} g(x) = \infty$ \qquad
(C) $\displaystyle\lim_{x\to\infty} g(x) = -\infty$ \qquad
(D) $\displaystyle\lim_{x\to-\infty} g(x) = -\infty$

\vspace{6pt}
\qn{25} The left end behavior of $g$: \quad \textit{Letter:}~\gradedblank[id=p25, type=choice, points=2, answer=B]{1.1cm}{13pt}
\hfill
\textbf{26.}\ The right end behavior of $g$: \quad \textit{Letter:}~\gradedblank[id=p26, type=choice, points=2, answer=A]{1.1cm}{13pt}

\vspace{10pt}
\qn{27} Complete the sentence: the
\gradedblank[id=p27a, type=short, points=1, answer=minimum, check=text, display={minimum}]{2.4cm}{1.1cm}
(\emph{maximum} or \emph{minimum}) value of $g$ is
\gradedblank[id=p27b, type=short, points=1, answer=-6, check=numeric]{1.4cm}{1.1cm},
and it occurs at $x = $
\gradedblank[id=p27c, type=short, points=1, answer=1, check=numeric]{1.4cm}{1.1cm}.

\end{document}
